\scanpage{108}\mspage{4}
en corr. $1-1$ avec les éléments de $(t_E^{\otimes2})^*$.

\uline{Corollaire 3} On a
\[
 \begin{aligned}
 F_\infty&\simeq k[P,P']\,/\,{P'}^2-(P^3+a_1P^2+a_2P^{\deleted{\unsure{2}}}+a_3)\\
 &\simeq k[P][P']\,/\,\bigl({P'}^2-(P^3+a_1P^2+a_2P^{\unsure{1|2}}+a_3)\bigr)
 \end{aligned} \tag*{(31)}
\]
l'isom étant donné par $P\mathrel{\unsure{-|\to}}p$, $P'\mathrel{\unsure{-|\to}}p'$. L'autom $f\to\check f$ correspond à l'autom. $P'\mapsto-P'$, $P\mapsto P$ de $k[P,P^{\unsure{+|'}}]/\ldots$ . La dérivation $\Theta_\delta$ est déduite par passage aux quotients de l'\deleted{application} dérivation
\[
 \Theta_{\deleted{\unsure{\delta}}}P=P',\qquad
 \Theta\deleted{\unsure{p}}P'=\frac12(3P^2+2a_1P+a_2) \tag*{(32)}
\]
sur $k[P,P']$, qui applique l'idéal $ {P'}^2-(P^3+a_1P^2+a_2P+a_3)$ en lui-$=$.

La formule (32) s'obtient en dérivant $ {p'}^2=p^3+a_1p^2+a_2p+a_3$ par $\Theta_\delta$, \deleted{et en divisant} \inserted{d'où}
\[
 2p'\Theta_\delta(p')=(3p^2+2a_1p+a_2)p'
\]
et en divisant par $p'$ la relation obtenue, d'où
\[
 p'=\tfrac12(3p^2+2a_1p+a_2)\,. \tag*{(33)}
\]
\deleted{\uline{Corollaire}}
\uline{Remarque} En termes des bases explicites plus haut des $F_n^+$ et $F_n^-$ $(25^1)$ donc des $F_n=F_n^+\oplus F_n^-$, on voit maintenant comment expliciter la loi de multiplication \inserted{en termes de $a_1,a_2,a_3$,} ainsi que
% archive sheet number: 107
